\chapter{Contexte et fondements}
\label{chapitre:contexte}

La détection automatisée suppose de connaître d'abord ce qu'elle cherche à repérer. L'enjeu
économique et zootechnique de la boiterie ouvre le chapitre, suivi de son étiologie et de ses
facteurs de risque. La notation locomotrice, méthode de référence du repérage, est ensuite
présentée avec ses limites à l'échelle d'un troupeau, avant les capteurs et les variables
comportementales mobilisés.

% =====================================================================
\section{Enjeu économique et zootechnique de la boiterie}
% =====================================================================

La boiterie bovine est un problème de santé majeur en élevage laitier, avec des conséquences
directes sur le bien-être animal, la productivité et la rentabilité des exploitations
\citep{bicalho2009lameness}.
Des études épidémiologiques rapportent des prévalences de l'ordre de 20\,\%. Celles-ci varient
selon les régions, les systèmes d'élevage et les critères de diagnostic utilisés
\citep{whay2003assessment,cook2003prevalence}. Les pathologies du pied, qu'elles soient
infectieuses (dermatite digitale, phlegmon interdigité) ou non infectieuses (ulcère de la
sole, maladie de la ligne blanche),
comptent parmi les principales causes de la boiterie chez les vaches laitières
\citep{merck_hoof_origin_lameness_cattle,flower2006effect}.

L'impact économique de la boiterie est substantiel et multidimensionnel. Le
Tableau~\ref{tab:impact_economique} rassemble les principales composantes documentées du
coût de la boiterie, d'après les estimations publiées dans la littérature.

\begin{table}[H]
    \centering
    \caption{Composantes documentées du coût de la boiterie.}
    \label{tab:impact_economique}
    \small
    \begin{tabular}{lcc}
        \toprule
        \textbf{Composante} & \textbf{Estimation publiée} & \textbf{Source} \\
        \midrule
        Perte de production laitière & $\approx$360 kg/lactation (305 j) & \cite{green2002financial} \\
        Trouble du pied clinique & 95~\$~US / cas & \cite{bruijnis2010assessing} \\
        Trouble du pied subclinique & 18~\$~US / cas & \cite{bruijnis2010assessing} \\
        Coût par vache et par an & 75~\$~US & \cite{bruijnis2010assessing} \\
        Ferme type de 65 vaches & 4\,899~\$~US / an & \cite{bruijnis2010assessing} \\
        \bottomrule
    \end{tabular}
\end{table}
\vspace{-0.9\baselineskip}

Le coût d'un épisode se répartit sur plusieurs éléments. La perte de production laitière est
la mieux documentée~: \citet{green2002financial} estiment une réduction moyenne d'environ
360~kg par lactation de 305~jours chez les vaches cliniquement boiteuses, la baisse
s'observant depuis plusieurs mois avant le diagnostic jusqu'à plusieurs mois après le
traitement. À cette perte s'ajoutent des coûts directs et indirects (soins, parage,
dégradation de la fertilité et réforme anticipée), dont l'ampleur varie selon la sévérité
et la chronicité de l'épisode. En intégrant ces postes dans un modèle de simulation,
\citet{bruijnis2010assessing} chiffrent le coût moyen d'un trouble du pied clinique à
95~\$~US par cas et celui d'un trouble subclinique à 18~\$~US, soit environ
75~\$~US par vache et par an et près de 4\,899~\$~US par an pour une
ferme type de 65~vaches.
Au-delà de l'impact financier, la boiterie soulève un enjeu éthique de bien-être animal~:
la douleur associée est reconnue comme l'une des principales sources de souffrance
chronique chez les vaches laitières \citep{whay2003assessment,broom2017animal}. Ces conséquences
justifient une compréhension fine des mécanismes de boiterie, afin d'orienter la
prévention, la détection et la prise en charge.

% =====================================================================
\section{Étiologie et facteurs de risque}
% =====================================================================

La boiterie a des causes multifactorielles. La Figure~\ref{fig:etiologie} les regroupe en
quatre catégories~: les causes directes infectieuses et non infectieuses ainsi que les facteurs
environnementaux et individuels. Les causes directes comprennent des affections bactériennes,
comme la \textbf{dermatite digitale} et le \textbf{phlegmon interdigité}, ainsi que des lésions
mécaniques, comme l'\textbf{ulcère de la sole} et la \textbf{maladie de la ligne blanche}.
D'autre part, l'exposition dépend de facteurs environnementaux comme le sol, l'hygiène et
l'humidité, tandis que la susceptibilité varie selon des facteurs individuels comme la parité,
l'âge et le stade de lactation. Cette diversité rend insuffisant le recours à un indicateur
unique et soutient une détection précoce fondée sur plusieurs signaux
\citep{merck_hoof_origin_lameness_cattle,dairynz_preventing_managing_lameness_guide_2025}.

\begin{figure}[H]
    \centering
    \resizebox{0.85\textwidth}{!}{%
    \begin{tikzpicture}[
        cat/.style={rectangle, draw=black!70, rounded corners=3pt, minimum height=0.9cm,
                    minimum width=3.2cm, align=center, font=\small, inner sep=10pt},
        factor/.style={rectangle, draw=black!40, rounded corners=2pt, minimum height=0.8cm,
                      minimum width=2.7cm, align=center, font=\small, inner sep=6pt},
        arrow/.style={-{Stealth[length=2mm]}, thick, black!60},
    ]
        % Titre central
        \node[cat, fill=rawred!15, minimum width=4cm, font=\small\bfseries] (boiterie)
            {Boiterie bovine};

        % Catégories
        \node[cat, fill=ifblue!12, above left=1.5cm and 1cm of boiterie] (infect)
            {Pathologies\\[-7pt]infectieuses};
        \node[cat, fill=loforange!12, above right=1.5cm and 1cm of boiterie] (noninfect)
            {Pathologies\\[-7pt]non infectieuses};
        \node[cat, fill=rulegreen!12, below left=1.5cm and 1cm of boiterie] (enviro)
            {Facteurs\\[-7pt]environnementaux};
        \node[cat, fill=profgray!12, below right=1.5cm and 1cm of boiterie] (indiv)
            {Facteurs\\[-7pt]individuels};

        % Liens
        \draw[arrow] (infect) -- (boiterie);
        \draw[arrow] (noninfect) -- (boiterie);
        \draw[arrow] (enviro) -- (boiterie);
        \draw[arrow] (indiv) -- (boiterie);

        % Sous-facteurs infectieux
        \node[factor, fill=ifblue!6, above=0.9cm of infect, xshift=-2.0cm] (dd)
            {Dermatite digitale};
        \node[factor, fill=ifblue!6, above=0.9cm of infect, xshift=2.0cm] (phleg)
            {Phlegmon interdigité};
        \draw[black!30] (dd) -- (infect);
        \draw[black!30] (phleg) -- (infect);

        % Sous-facteurs non infectieux
        \node[factor, fill=loforange!6, above=0.9cm of noninfect, xshift=-2.0cm] (ulcere)
            {Ulcère de la sole};
        \node[factor, fill=loforange!6, above=0.9cm of noninfect, xshift=2.0cm] (ligne)
            {Maladie ligne blanche};
        \draw[black!30] (ulcere) -- (noninfect);
        \draw[black!30] (ligne) -- (noninfect);

        % Sous-facteurs environnementaux
        \node[factor, fill=rulegreen!6, below=0.9cm of enviro, xshift=-2.0cm] (sol)
            {Type de sol};
        \node[factor, fill=rulegreen!6, below=0.9cm of enviro, xshift=2.0cm] (hygiene)
            {Hygiène, humidité};
        \draw[black!30] (sol) -- (enviro);
        \draw[black!30] (hygiene) -- (enviro);

        % Sous-facteurs individuels
        \node[factor, fill=profgray!6, below=0.9cm of indiv, xshift=-2.0cm] (parite)
            {Parité, âge};
        \node[factor, fill=profgray!6, below=0.9cm of indiv, xshift=2.0cm] (lactation)
            {Stade de lactation};
        \draw[black!30] (parite) -- (indiv);
        \draw[black!30] (lactation) -- (indiv);
    \end{tikzpicture}%
    }
    \caption{Étiologie multifactorielle de la boiterie bovine.}
    \label{fig:etiologie}
\end{figure}

% =====================================================================
\section{Notation locomotrice et détection conventionnelle}
% =====================================================================

La notation locomotrice reste la référence pratique du repérage de la boiterie. Cette
section précise d'abord pourquoi un repérage précoce importe, puis ce que l'observation
humaine permet et ce qu'elle ne permet pas de maintenir à l'échelle d'un troupeau.

\subsection{Importance de la détection précoce}
La littérature sur la boiterie bovine souligne l'importance d'une détection précoce
\citep{oleary2020accelerometers,qiao2021review} pour
limiter la dégradation du bien-être animal et les impacts de production. Les ressources
vétérinaires de référence et les guides techniques convergent sur deux idées:
une surveillance régulière est indispensable, et la prise en charge doit être
structurée et rapide selon la sévérité du cas
\citep{merck_lameness_overview,dairynz_treating_lameness}.

Dans la grille de notation locomotrice, le score~2 correspond à une boiterie légère et constitue
le seuil d'un repérage précoce. Une intervention à ce stade peut éviter que les cas ne deviennent
chroniques. Les études montrent qu'un
retard de traitement de quelques jours peut multiplier la durée de récupération et le
coût associé \citep{green2002financial}. Les capteurs portés et les systèmes de mesure de
la démarche permettent de quantifier objectivement des modifications d'activité ou de
locomotion associées à la boiterie
\citep{chapinal2010measurement,pastell2009measurement}. Le délai avec lequel ces
modifications deviennent détectables avant une observation clinique dépend toutefois du
capteur, de la définition du cas et du protocole; il ne peut donc pas être généralisé à
partir de ces travaux \citep{oleary2020accelerometers,alsaaod2019automatic}. En pratique, les boiteries légères restent difficiles à repérer de
façon régulière par la seule observation visuelle.

Des revues de synthèse montrent aussi que la boiterie figure parmi les cas d'usage les plus
importants en élevage de précision, mais aussi parmi les plus difficiles à stabiliser, car
les signaux disponibles restent indirects, multifactoriels et fortement dépendants du
contexte d'élevage \citep{palma2025scoping,michelena2025plf_review}.

\subsection{Observation et notation locomotrice}
Les approches conventionnelles reposent principalement sur l'observation humaine et les
grilles de notation locomotrice. Une grille à cinq niveaux largement utilisée a été proposée
pour standardiser cette observation \citep{sprecher1997lameness}. Elle attribue un score allant
de 1, pour une locomotion normale, à 5, pour une boiterie sévère. Cette notation tient
principalement compte de la courbure du dos à l'arrêt et en mouvement ainsi que de
l'irrégularité de la démarche \citep{flower2006gait}. La fiche opérationnelle de
l'\textit{Agriculture and Horticulture Development Board} (AHDB) propose une grille
complémentaire adaptée au contexte britannique et associe le score~2 à une
intervention rapide, généralement dans les 48~heures
\citep{ahdb_mobility_score_sheet_48h_score2}.

Les approches conventionnelles restent pertinentes en élevage, car elles sont simples à
déployer et directement interprétables par les équipes de terrain. Elles présentent toutefois
plusieurs limites opérationnelles. La première contrainte tient aux différences d'appréciation
entre observateurs. Même avec des grilles standardisées, la concordance entre évaluateurs peut varier significativement
\citep{flower2006effect,march2007interobserver}. La deuxième concerne le temps requis:
dans les grands troupeaux, l'observation systématique de chaque vache à chaque traite est
difficile à maintenir. Enfin, les boiteries légères (score~2) restent difficiles à repérer
de manière reproductible par la seule observation visuelle
\citep{whay2003assessment,march2007interobserver}.

Pour pallier ces limites, le recours aux signaux issus de capteurs s'est avéré nécessaire afin de
compléter la surveillance humaine sans la remplacer.

\subsection{Lecture opérationnelle du score locomoteur}
Au-delà de sa fonction descriptive, l'évaluation locomotrice est aussi un outil de
priorisation. Dans la fiche opérationnelle de l'AHDB, le passage du score~1
(normal) au score~2 (boiterie légère) appelle une intervention rapide,
généralement dans les 48~heures
\citep{ahdb_mobility_score_sheet_48h_score2}. Cette logique explique
pourquoi un système automatisé peut viser une fonction plus prudente que la reproduction
d'une note clinique: signaler une dégradation comportementale suffisamment tôt pour
déclencher une vérification sur le terrain.

La Figure~\ref{fig:sprecher_scale} présente la lecture opérationnelle du score
locomoteur: le passage vers le score~2 constitue la zone d'intérêt pour une alerte
précoce, car il marque l'entrée dans une situation où l'observation humaine ou l'examen
du pied devient prioritaire
\citep{sprecher1997lameness,ahdb_mobility_score_sheet_48h_score2}.

\begin{figure}[H]
    \centering
    \begin{tikzpicture}[
        score/.style={rectangle, draw=black!60, rounded corners=3pt, minimum width=2.45cm,
                      minimum height=1.7cm, align=center, inner sep=8pt, font=\small},
        arrow/.style={-{Stealth[length=2.5mm]}, thick, black!55}
    ]
        \node[score, fill=rulegreen!12] (s1) at (0,0.2) {\textbf{Score 1}\\[2pt]Normal};
        \node[score, fill=ifblue!10] (s2) at (3.2,0.2) {\textbf{Score 2}\\[2pt]Léger};
        \node[score, fill=loforange!14] (s3) at (6.4,0.2) {\textbf{Score 3}\\[2pt]Modéré};
        \node[score, fill=rawred!14] (s4) at (9.6,0.2) {\textbf{Score 4}\\[2pt]Marqué};
        \node[score, fill=rawred!25] (s5) at (12.8,0.2) {\textbf{Score 5}\\[2pt]Sévère};

        \draw[arrow] (1.35,0.2) -- (1.95,0.2);
        \draw[arrow] (4.55,0.2) -- (5.15,0.2);
        \draw[arrow] (7.75,0.2) -- (8.35,0.2);
        \draw[arrow] (10.95,0.2) -- (11.55,0.2);

        \draw[decorate, decoration={brace, amplitude=6pt}, thick, ifblue]
            (1.7,1.7) -- (4.7,1.7)
            node[midway, above=8pt, align=center, font=\footnotesize\bfseries]
            {Zone d'intérêt pour\\la détection précoce};
    \end{tikzpicture}

    \caption[Seuil opérationnel du score locomoteur]{Seuil opérationnel du score
    locomoteur pour la détection précoce.}
    \label{fig:sprecher_scale}
\end{figure}

Le Tableau~\ref{tab:score_operationnel} précise la lecture clinique simplifiée et
l'action associée à chaque score, d'après la grille de \citet{sprecher1997lameness} et
les recommandations AHDB \citep{ahdb_mobility_scoring_dairy_cows}.

\begin{table}[H]
    \centering
    \caption{Lecture clinique et opérationnelle simplifiée du score locomoteur.}
    \label{tab:score_operationnel}
    \small
    \begin{tabular}{>{\raggedright\arraybackslash}p{1.4cm}
                    >{\raggedright\arraybackslash}p{6.0cm}
                    >{\raggedright\arraybackslash}p{5.2cm}}
        \toprule
        \textbf{Score} & \textbf{Lecture clinique simplifiée} & \textbf{Lecture opérationnelle} \\
        \midrule
        1 & Dos plat en station et en marche. & Surveillance usuelle. \\
        2 & Dos plat à l'arrêt, légèrement arqué en marche. & Intervention rapide: zone clé pour la détection précoce. \\
        3 & Dos arqué en permanence. & Vérification prioritaire. \\
        4 & Atteinte locomotrice visible. & Prise en charge urgente. \\
        5 & Appui très limité. & Prise en charge immédiate. \\
        \bottomrule
    \end{tabular}
\end{table}

La lecture par seuil transforme la notation locomotrice en problème d'alerte: il ne
s'agit plus seulement de décrire une boiterie, mais de repérer l'entrée dans une zone
d'intervention. Ce déplacement justifie le recours à des signaux capteurs et à des
méthodes d'apprentissage capables de suivre les changements comportementaux dans le
temps.

% =====================================================================
\section{Capteurs et variables comportementales}
% =====================================================================

\label{sec:technologies_capteurs}

Cette section situe brièvement les principales familles de capteurs utilisées pour la détection
de la boiterie. Elle décrit ensuite le capteur retenu dans ce projet, IceQube, ainsi que les
variables comportementales qu'il produit.

Les systèmes de surveillance comportementale reposent sur plusieurs familles de capteurs,
dont les caractéristiques conditionnent directement les variables disponibles et les
performances atteignables. Le Tableau~\ref{tab:technologies_capteurs} présente les principales
technologies et leurs caractéristiques opérationnelles. On distingue trois grandes familles:
les \textbf{capteurs portés par l'animal}, tels que les accéléromètres, les podomètres et les
colliers; les \textbf{capteurs intégrés à l'infrastructure}, comme les tapis de pesée et les
systèmes de traite; et les \textbf{systèmes sans contact}, comme les caméras et la vision par
ordinateur. Les accéléromètres
portés constituent la technologie la plus étudiée dans la littérature récente
\citep{oleary2020accelerometers,alsaaod2019automatic}.

\begin{table}[H]
    \centering
    \small
    \caption{Principales familles de capteurs étudiées.}
    \label{tab:technologies_capteurs}
    \begin{tabular}{>{\raggedright\arraybackslash}p{2.8cm}
                    >{\raggedright\arraybackslash}p{3.2cm}
                    >{\raggedright\arraybackslash}p{3.2cm}
                    >{\raggedright\arraybackslash}p{2.6cm}
                    >{\raggedright\arraybackslash}p{2.6cm}}
        \toprule
        \textbf{Technologie} & \textbf{Variables mesurées} & \textbf{Avantages} &
        \textbf{Limites} & \textbf{Études représentatives} \\
        \midrule
        Accéléromètre (patte/collier)
        & Activité, pas, transitions couché/debout, durée couchage
        & Continu, individuel, coût modéré
        & Dérive, position variable
        & \cite{alsaaod2012electronic,balasso2023behaviour} \\
        \addlinespace
        Podomètre
        & Nombre de pas, distance parcourue
        & Simple, robuste
        & Signal limité, faible résolution
        & \cite{flower2006effect,mazrier2006reliability} \\
        \addlinespace
        Passerelle sensible à la pression
        & Pressions plantaires et paramètres spatio-temporels de la démarche
        & Objectif, quantitatif
        & Infrastructure fixe, coût élevé
        & \cite{maertens2011gaitwise} \\
        \addlinespace
        Vision par ordinateur
        & Posture, arc du dos, démarche
        & Sans contact, riche en information
        & Sensible aux conditions, calibration
        & \cite{qiao2021review} \\
        \addlinespace
        Système automatisé de traite
        & Production laitière, conductivité, fréquentation
        & Déjà installé, données continues
        & Signal indirect pour la boiterie
        & \cite{lavrova2023lameness,vanhertem2013lameness} \\
        \bottomrule
    \end{tabular}
\end{table}

Dans le cadre de ce mémoire, les données proviennent de capteurs IceQube
(IceRobotics Ltd, Édimbourg, Royaume-Uni), à base d'accéléromètres tri-axiaux conçus
spécifiquement pour le suivi comportemental des vaches laitières. La
Figure~\ref{fig:iceqube_capteur} montre la représentation schématique du capteur IceQube,
dont le boîtier étanche contient l'accéléromètre tri-axial et se fixe à la patte du bovin par un
bracelet en plastique muni d'une boucle et de trous d'ajustement. La
Figure~\ref{fig:iceqube_position} illustre le positionnement sur le métatarse de la patte
arrière~; les axes $x$, $y$, $z$ de l'accéléromètre permettent de distinguer la position
debout (patte verticale, $\vec{g}$ aligné avec $z$) de la position couchée (patte
horizontale).

\begin{figure}[H]
    \centering
    \resizebox{0.7\textwidth}{!}{%
    \begin{tikzpicture}[>=Stealth]
        % --- Bracelet gauche (vers la boucle) ---
        \fill[ifblue!35, rounded corners=2pt]
            (-3.0,0.65) rectangle (0,1.35);
        % Boucle de fermeture
        \fill[ifblue!50, rounded corners=2pt]
            (-3.6,0.5) rectangle (-3.0,1.5);
        \draw[ifblue!70, thick, rounded corners=1pt]
            (-3.45,0.7) rectangle (-3.15,1.3);

        % --- Bracelet droit (trous) ---
        \fill[ifblue!35, rounded corners=2pt]
            (4.0,0.65) rectangle (7.0,1.35);
        % Trous d'ajustement (ovales)
        \foreach \x in {4.6, 5.2, 5.8, 6.4} {
            \fill[white] (\x,1.0) ellipse (0.12 and 0.15);
        }

        % --- Boîtier central (vue de face) ---
        \fill[ifblue!50, rounded corners=4pt]
            (0,0) rectangle (4.0,2.0);
        % Bord supérieur légèrement plus clair
        \fill[ifblue!38, rounded corners=3pt]
            (0.1,1.55) rectangle (3.9,1.95);

        % Étiquette blanche
        \fill[white, rounded corners=2pt]
            (0.6,0.35) rectangle (3.4,1.4);
        \node[font=\small\bfseries, black] at (2.0,1.05) {IceQube};
        \node[font=\footnotesize, black!60] at (2.0,0.7) {IceRobotics};

        % --- Annotations ---
        \draw[->, black!55, thick]
            (4.5,2.5) -- (3.5,2.05)
            node[pos=0, anchor=south west, font=\footnotesize, black!70]
            {Boîtier étanche};

        \draw[->, black!55, thick]
            (-3.6,0.0) -- (-3.45,0.45)
            node[pos=0, anchor=north, font=\footnotesize, black!70]
            {Boucle};

        \draw[->, black!55, thick]
            (7.0,0.0) -- (5.8,0.6)
            node[pos=0, anchor=north, font=\footnotesize, black!70]
            {Trous d'ajustement};
    \end{tikzpicture}%
    }%
    \caption{Schéma du capteur IceQube.}
    \label{fig:iceqube_capteur}
\end{figure}

\begin{figure}[H]
    \centering
    \resizebox{0.48\textwidth}{!}{%
    \begin{tikzpicture}[>=Stealth]
        % --- Patte (forme simple, allongée verticalement) ---
        \fill[brown!12, rounded corners=5pt]
            (0,0) -- (0,6.0) -- (0.4,6.5) -- (1.4,6.5) -- (1.8,6.0)
            -- (1.8,0) -- (1.5,-0.4) -- (0.3,-0.4) -- cycle;
        \draw[brown!30, thick, rounded corners=5pt]
            (0,0) -- (0,6.0) -- (0.4,6.5) -- (1.4,6.5) -- (1.8,6.0)
            -- (1.8,0) -- (1.5,-0.4) -- (0.3,-0.4) -- cycle;

        % Jarret (ligne courte, texte à droite)
        \draw[brown!50, thick] (-0.2,4.5) -- (2.2,4.5);
        \node[font=\footnotesize, brown!80!black, anchor=west] at (2.4,4.5) {Jarret};

        % Zone métatarse (lignes courtes à droite seulement)
        \draw[brown!45, dashed] (1.9,1.0) -- (2.8,1.0);
        \draw[brown!45, dashed] (1.9,3.2) -- (2.8,3.2);
        \draw[decorate, decoration={brace, amplitude=4pt, mirror}, brown!60]
            (2.9,1.0) -- (2.9,3.2)
            node[midway, right=5pt, font=\footnotesize, brown!80!black] {Métatarse};

        % --- Bracelet autour de la patte ---
        \fill[ifblue!35, rounded corners=2pt]
            (-0.6,1.7) rectangle (0,2.5);
        \fill[ifblue!35, rounded corners=2pt]
            (1.8,1.7) rectangle (2.4,2.5);
        % Boucle côté gauche
        \fill[ifblue!50, rounded corners=1pt]
            (-0.8,1.8) rectangle (-0.6,2.4);

        % --- Capteur (sur la face avant de la patte) ---
        \fill[ifblue!50, rounded corners=3pt]
            (0.0,1.6) rectangle (1.8,2.6);
        \fill[white, rounded corners=1pt]
            (0.25,1.75) rectangle (1.55,2.45);
        \node[font=\footnotesize\bfseries, black] at (0.9,2.1) {IceQube};

        % --- Axes de l'accéléromètre (en bas à gauche, bien isolés) ---
        \coordinate (AX) at (-2.2,0.0);

        % z vers le haut
        \draw[->, rawred, very thick] (AX) -- ++(0,1.4)
            node[above, font=\small, rawred] {$z$};

        % x vers la droite
        \draw[->, rulegreen, very thick] (AX) -- ++(1.2,0)
            node[right, font=\small, rulegreen] {$x$};

        % y vers l'avant (hors plan)
        \draw[->, ifblue!80!black, very thick] (AX) -- ++(-0.7,-0.7)
            node[below left, font=\small, ifblue!80!black] {$y$};

        % --- Gravité (à droite, au-dessus du jarret) ---
        \draw[->, black!65, thick, dashed] (3.8,6.0) -- ++(0,-1.3)
            node[midway, right=2pt, font=\footnotesize, black!75] {$\vec{g}$};
    \end{tikzpicture}%
    }%
    \caption{Positionnement du capteur IceQube.}
    \label{fig:iceqube_position}
\end{figure}

L'accéléromètre interne mesure l'accélération sur trois axes. La composante gravitationnelle
distingue les positions couchée et debout, tandis que les variations d'accélération fournissent
les indicateurs d'activité. L'algorithme propriétaire agrège les mesures brutes; le présent
projet exploite des intervalles de 15~minutes. Cette résolution, justifiée à la
Section~\ref{sec:resolution_temporelle}, est un compromis analytique plutôt qu'une fréquence
clinique démontrée. Des études de capteurs portés rapportent une forte concordance avec
l'observation directe du comportement couché et de la posture couché/debout
\citep{ledgerwood2010lying,alsaaod2012electronic}, ainsi que des différences de couchage
associées à la boiterie \citep{blackie2011lying}.

Chaque intervalle de 15~minutes génère sept colonnes, dont cinq sont utilisées dans la chaîne de
traitement: \textit{Steps}, le nombre de pas détectés; \textit{Motion Index}, un indice
synthétique d'activité locomotrice dérivé des variations d'accélération (unité propriétaire
IceRobotics); \textit{Lying Time} et \textit{Standing Time}, les durées cumulées en position
couchée et debout, chacune bornée à 15~minutes; et \textit{Transitions}, le nombre total de
changements de posture couché/debout.
Chez des veaux sevrés équipés d'un accéléromètre IceTag, le \textit{Motion Index} a permis
d'identifier le jeu au moyen de seuils validés \citep{mckay2024play}. Ce résultat confirme
la valeur discriminante du signal, mais aussi son absence de spécificité pour la boiterie.
Les deux colonnes supplémentaires du fichier IceQube (\textit{Transitions Up} et
\textit{Transitions Down}) détaillent la direction des transitions; la chaîne de traitement
utilise plutôt leur somme.
Les données agrégées, stockées dans la mémoire interne du capteur, sont récupérées périodiquement
par CowAlert (IceRobotics), puis exportables depuis son interface web au format tabulaire (CSV),
directement exploitable par la chaîne de traitement.

Le modèle IceTag, prédécesseur de l'IceQube, fournit des variables similaires (pas, temps
couché, transitions) et a été utilisé dans plusieurs travaux sur le comportement et la
locomotion bovins \citep{beer2016objective,alsaaod2012electronic}. Le choix de l'IceQube
est d'abord déterminé par la disponibilité du corpus et s'accorde avec la place importante
des accéléromètres portés dans la surveillance locomotrice
\citep{oleary2020accelerometers,bewley2015precision}. Les vaches boiteuses modifient notamment
leur comportement de couchage \citep{ito2010lying}, tandis que des lésions de la corne sont
associées à des changements comportementaux mesurables \citep{nechanitzky2016analysis}.
Ces observations justifient l'exploitation des variables IceQube comme indicateurs indirects
de changements locomoteurs, sans établir à elles seules un délai universel de détection
préclinique ni une spécificité clinique. Le chapitre suivant examine les approches proposées
pour transformer ces signaux en alertes, ainsi que leurs performances et leurs limites
méthodologiques.
